\section{Asking a model for an independent check}
\label{sec:automation}

The pipeline was written by people. We tested one part of the question for
automated authors: asked to add an independent second path for the disposition
count, does a code-generating model produce one whose meaning comes from
somewhere other than the path it was shown? Reuse is the expected answer, since
not duplicating a definition is correct practice by every other standard, and
generated test oracles are known to encode the behaviour of the code in front of
them rather than the behaviour intended \citep{konstantinou2024oracles}.

\paragraph{A preliminary run, and what was wrong with it.}
A preliminary run of 30 trials on 19~August used a coding-agent CLI with the
aliases \texttt{sonnet} and \texttt{haiku} and labelled outputs by pattern
matching over the generated source. It had four problems. The aliases were not resolved to model
identifiers in any record we kept. The CLI's default tools were not disabled, so
we cannot show that no trial read files. The prompt paraphrased the consumer
function: it counted only explicit membership of the gone and in-orbit sets,
where the real code treated every other status as in orbit. On the frozen data
that paraphrase returns 208, not 932, so the trials never saw the defect they
were scored against. And the loader bodies were elided, so no trial could see
GCAT's column names. What survives is narrower than we claimed. None of the 30
outputs introduced a status code absent from the prompt, and executed on the
real files, 22 of them return exactly the 208 of the path they were shown,
including all 19 that ran under the neutral and ``independent'' requests. When
code sharing was forbidden, three of ten reproduced 208 and seven diverged, none
by consulting the source: four guessed a column name and return 0, one parses
differently and returns 189, one crashes, and one reads GCAT's phase end date as
a decay date and returns 795 (Appendix~\ref{app:experiment}).

\paragraph{A controlled replication.}
We reran the experiment with the pre-correction code verbatim, loaders
included, so that the path shown returns the published 932. Each model was
pinned by identifier and called through the same CLI with tools disabled, a
one-line system prompt and an empty working directory: \texttt{claude-opus-5-5},
\texttt{claude-sonnet-5} and \texttt{claude-haiku-4-5-20251001}. Beside the three original requests,
two new conditions append GCAT's own status definitions verbatim, which is the
remedy this paper recommends. Outcomes are behavioural. Each generated path is
executed on a synthetic catalogue in which each pairing of a status code with a
CelesTrak state holds a distinct power of two of objects, so any count decodes to
exactly which pairings the path counted, and then on the frozen real files
(Appendix~\ref{app:experiment}).

\begin{table}[t]
\centering
\small
\caption{Controlled replication. Each cell is the number of generated paths that
return the published, defective 932 on the frozen data, out of the trials run.
``+ definitions'' appends GCAT's status definitions verbatim.
All 72 paths that return 932 count the same status pairings as the path they
were shown, on the probe or, for one, on the real files.}
\label{tab:reuse}
\begin{tabular}{lccc}
\toprule
Request & \texttt{opus-5-5} & \texttt{sonnet-5} & \texttt{haiku-4-5} \\
\midrule
Neutral                                  & 5/5 & 5/5 & 5/5 \\
``Independent'' path                     & 5/5 & 5/5 & 5/5 \\
Sharing forbidden                        & 4/5 & 4/5 & 5/5 \\
``Independent'' + definitions & 5/5 & 5/5 & 5/5 \\
Sharing forbidden + definitions & 4/5 & 5/5 & 5/5 \\
\bottomrule
\end{tabular}
\end{table}

Table~\ref{tab:reuse} gives the result. Of 75 trials, 72 return 932. Of the
other three, one \texttt{claude-sonnet-5} path returns 0 because its
hand-written file parser fails, and two \texttt{claude-opus-5-5} paths refuse
to produce a count. Only one of the 75 would have exposed the defect: with code
sharing forbidden, an Opus path wrote its own completeness check, and on the
real files it stops and lists the unclassified codes, among them \texttt{E},
\texttt{C}, \texttt{DK} and \texttt{ATT}. That is the first gate of
Section~\ref{sec:gates}, produced unprompted, once.

The definitions did not change what was computed. Of the 30 paths written with
GCAT's definitions in the prompt, 29 return 932, and every one of the 30 counts
\texttt{E}, \texttt{C}, \texttt{DK}, \texttt{ATT}, \texttt{TFR} and
\texttt{GRP} exactly as the shown path does wherever the probe can see it. The
paths read the documentation and kept the comparison. One Opus path annotated
\texttt{E} as ``exploded; the object continues as a debris fragment'' and
\texttt{C} as ``collided; may or may not survive, so not treated as gone'',
both correct, then placed both, with the docking codes, in a class of objects
``still in space or transition'' and compared that class with CelesTrak's decay
flag as the defective path does. On the real files its own completeness check
stops it on the undocumented suffixes \texttt{L?} and \texttt{R?}; on the
probe, which has none, it counts exactly what the defective path counts. A \texttt{claude-sonnet-5} path's docstring says
it ``re-derives the GCAT status taxonomy straight from the phase
documentation''; it returns 932.

\paragraph{What this does and does not establish.}
On this task, asking for an independent path, forbidding shared code, and
supplying the source's definitions each produced paths that, where they worked,
computed what the defective path computes, while looking to a code reviewer like
separate implementations with new names and comments citing the documentation.
It does not establish a rate for models in general (Section~\ref{sec:limits}). If a check is
meant to be independent, independence has to be specified as a property of its
inputs and verified by what the check computes, not requested; models also
struggle to correct their own reasoning without external feedback
\citep{huang2024cannotselfcorrect}.
